% 分类目录：ml
\documentclass[12pt, a4paper]{article}
\usepackage[margin=2.5cm]{geometry}
\usepackage{ctex}
\usepackage{amsmath, amssymb}
\usepackage{listings}
\usepackage{xcolor}
\usepackage{tabularx}
\usepackage{hyperref}

\lstset{
    language=Python,
    basicstyle=\ttfamily\small,
    keywordstyle=\color{blue},
    commentstyle=\color{green!50!black},
    stringstyle=\color{red},
    showstringspaces=false,
    breaklines=true,
    numbers=left,
    numberstyle=\tiny\color{gray},
    frame=single
}

\title{知识点：K-近邻算法 (K-Nearest Neighbors)}
\author{}
\date{}

\begin{document}
\maketitle

\section{核心定义}
\textbf{中文定义}：K-近邻算法（KNN）是一种基于距离度量的监督学习算法。它的核心思想是“物以类聚”：即对于一个待分类的样本，在训练集中寻找与之最接近的 $k$ 个样本，根据这 $k$ 个样本的类别进行投票或加权平均，从而预测待分类样本的类别或数值。

\textbf{英文定义}：K-Nearest Neighbors (KNN) is a distance-based supervised learning algorithm. Its central idea is "birds of a feather flock together": for a given sample, the algorithm finds the $k$ most similar samples (neighbors) in the training set and predicts the label or value based on the majority vote or weighted average of these neighbors.

\section{三大核心要素}

\subsection{距离度量 (Distance Metric)}
距离的选择直接影响模型性能。常用的距离包括：
\begin{itemize}
    \item \textbf{欧元距离 (Euclidean Distance)}：
    $$ L_2(x_i, x_j) = \sqrt{\sum_{l=1}^d (x_i^{(l)} - x_j^{(l)})^2} $$
    \item \textbf{曼哈顿距离 (Manhattan Distance)}：
    $$ L_1(x_i, x_j) = \sum_{l=1}^d |x_i^{(l)} - x_j^{(l)}| $$
    \item \textbf{闵可夫斯基距离 (Minkowski Distance)}：
    $$ L_p(x_i, x_j) = \left( \sum_{l=1}^d |x_i^{(l)} - x_j^{(l)}|^p \right)^{1/p} $$
\end{itemize}

\subsection{$k$ 值的选择}
$k$ 值的大小决定了模型的复杂程度：
\begin{itemize}
    \item \textbf{较小的 $k$ 值}：模型较复杂，容易受到噪声干扰，导致\textbf{过拟合}（Variance 高）。
    \item \textbf{较大的 $k$ 值}：模型较简单，容易模糊类别边界，导致\textbf{欠拟合}（Bias 高）。
    \item \textbf{最佳实践}：通常通过交叉验证（Cross-validation）选取最优 $k$ 值。
\end{itemize}

\subsection{决策规则 (Decision Rule)}
\begin{itemize}
    \item \textbf{分类任务}：多数表决法（Majority Voting）。也可引入基于距离的反比加权。
    \item \textbf{回归任务}：平均值法（Averaging）。
\end{itemize}

\section{ lazy Learning (惰性学习)}
KNN 是一种典型的“惰性学习”算法。它在训练阶段几乎不做任何计算，只是简单地存储数据。真正的计算发生在\textbf{预测阶段}，这导致预测的计算量非常大。

\section{算法加速：KD 树 (K-Dimensional Tree)}
为了避免全样本暴力扫描（Brute-force），常用 KD 树进行搜索优化。
\begin{itemize}
    \item \textbf{构造}：递归地选择方差最大的维度作为切分轴，利用中位数构造平衡二叉树。
    \item \textbf{搜索}：通过“回溯”机制在树结构中寻找最近邻，大幅减少计算量。
\end{itemize}

\section{优缺点与关键技巧}

\subsection{优缺点分析}
\begin{table}[h]
\centering
\begin{tabularx}{\textwidth}{|l|X|}
\hline
\textbf{优点} & 简单直观、无需参数估计、对异常值（在 $k$ 较大时）有一定鲁棒性、适用于多分类任务。 \\
\hline
\textbf{缺点} & 计算复杂度及空间复杂度高；对特征缩放敏感；无法处理类别极其不平衡的数据。 \\
\hline
\end{tabularx}
\end{table}

\subsection{关键技巧：特征缩放 (Scaling)}
由于 KNN 基于距离，尺度大的特征会主导计算结果。因此，使用 KNN 前\textbf{必须}进行标准化或归一化。

\section{关键代码片段 (Python)}
\begin{lstlisting}
from sklearn.neighbors import KNeighborsClassifier
from sklearn.preprocessing import StandardScaler

# 1. 数据缩放 (极其重要)
scaler = StandardScaler()
X_train_scaled = scaler.fit_transform(X_train)

# 2. 初始化与训练
knn = KNeighborsClassifier(n_neighbors=5, metric='minkowski', p=2)
knn.fit(X_train_scaled, y_train)

# 3. 预测
y_pred = knn.predict(X_test_scaled)
\end{lstlisting}

\section{深度面试题}

\textbf{问：KNN 中的 $k$ 取 1 和 $k$ 取 $N$ 分别会发生什么？}\\
答：$k=1$ 时模型极度复杂，训练集误差为 0 但泛化能力差（强过拟合）；$k=N$ 时算法退化为“盲猜”，即预测结果永远是样本最多的那一类（强欠拟合）。

\textbf{问：KNN 在高维空间中表现如何？}\\
答：表现极差，这被称为“维度灾难”。在高维空间下，样本点之间的距离趋于相等，欧氏距离失效，且搜索效率呈指数级下降。

\section{参考文献}
\begin{enumerate}
    \item Cover, T., \& Hart, P. (1967). \textit{Nearest Neighbor Pattern Classification}. IEEE Transactions on Information Theory.
    \item Bentley, J. L. (1975). \textit{Multidimensional Binary Search Trees Used for Associative Searching}. Communications of the ACM.
    \item Altman, N. S. (1992). \textit{An Introduction to Kernel and Nearest-neighbor Nonparametric Regression}. The American Statistician.
\end{enumerate}

\end{document}
